\textbf{Phần 1. Trắc nghiệm lựa chọn}
\setcounterpageref{ex}{0}

\begin{ex} %[2H9N47]
Khối lập phương có bao nhiêu mặt?
\choice
{$4$ mặt}
{\True $6$ mặt vuông bằng nhau}
{$8$ mặt}
{$12$ mặt}
\loigiai{
Khối lập phương có 6 mặt vuông bằng nhau.\\
\textbf{Đáp án của câu hỏi là B}
}
\end{ex}

\begin{ex} %[2H9N47]
Vật nào sau đây lăn được về mọi phía?
\choice
{Hộp bánh hình chữ nhật}
{Khối rubik}
{\True Quả bóng bàn (khối cầu)}
{Lon sữa bò}
\loigiai{
Quả bóng bàn có dạng khối cầu lăn về mọi phía.\\
\textbf{Đáp án của câu hỏi là C}
}
\end{ex}

\begin{ex} %[2H9H47]
Khối trụ khi đặt nằm ngang thì có thể
\choice
{Không di chuyển}
{\True Lăn được}
{Xếp chồng vững chắc}
{Lăn về mọi hướng}
\loigiai{
Khi nằm ngang thì khối trụ lăn được.\\
\textbf{Đáp án của câu hỏi là B}
}
\end{ex}

\begin{ex} %[2H9H47]
Trong các khối: khối cầu, khối trụ, khối hộp chữ nhật, khối lập phương. Khối nào KHÔNG có mặt phẳng nào?
\choice
{\True Khối cầu}
{Khối trụ}
{Khối lập phương}
{Khối hộp chữ nhật}
\loigiai{
Khối cầu không có mặt phẳng nào.\\
\textbf{Đáp án của câu hỏi là A}
}
\end{ex}

\begin{ex} %[2H9H47]
Để xếp một ngôi nhà đồ chơi gồm phần thân là khối hộp chữ nhật và mái nhà là một khối tam giác, cần mấy khối hình?
\choice
{$1$ khối}
{$3$ khối}
{\True $2$ khối}
{$4$ khối}
\loigiai{
Cần 2 khối hình ghép lại.\\
\textbf{Đáp án của câu hỏi là C}
}
\end{ex}

\begin{ex} %[2H9H47]
Một hộp phấn có dạng hình khối nào?
\choice
{\True Khối hộp chữ nhật hoặc lập phương}
{Khối cầu}
{Khối trụ}
{Hình tròn}
\loigiai{
Hộp phấn có dạng khối hộp chữ nhật hoặc lập phương.\\
\textbf{Đáp án của câu hỏi là A}
}
\end{ex}

\begin{ex} %[2H9V47]
Bác thợ mộc có 8 khúc gỗ hình khối trụ. Bác cưa đôi mỗi khúc gỗ. Hỏi bác thu được bao nhiêu khúc gỗ hình khối trụ nhỏ?
\choice
{$14$ khúc}
{$18$ khúc}
{\True $16$ khúc}
{$12$ khúc}
\loigiai{
$8 \times 2 = 16$ khúc.\\
\textbf{Đáp án của câu hỏi là C}
}
\end{ex}

\begin{ex} %[2H9V47]
Trong một giỏ đồ chơi có 5 quả bóng, 4 hộp sữa bột hình trụ và 3 khối rubik. Tổng số khối hình trong giỏ là
\choice
{$11$ khối}
{$13$ khối}
{$14$ khối}
{\True $12$ khối}
\loigiai{
$5 + 4 + 3 = 12$ khối.\\
\textbf{Đáp án của câu hỏi là D}
}
\end{ex}

\begin{ex} %[2H9V47]
Một thùng có 20 lon nước ngọt xếp đều thành 4 hàng. Mỗi hàng có số lon nước ngọt là
\choice
{$4$ lon}
{\True $5$ lon}
{$6$ lon}
{$7$ lon}
\loigiai{
$20 : 4 = 5$ lon.\\
\textbf{Đáp án của câu hỏi là B}
}
\end{ex}

\begin{ex} %[2H9V47]
Muốn ghép 8 khối lập phương nhỏ cạnh $1\text{ cm}$ thành một khối lập phương lớn hơn, cạnh của khối lập phương lớn là
\choice
{$1\text{ cm}$}
{$3\text{ cm}$}
{$4\text{ cm}$}
{\True $2\text{ cm}$}
\loigiai{
Cạnh khối lập phương lớn: $1 + 1 = 2\text{ cm}$.\\
\textbf{Đáp án của câu hỏi là D}
}
\end{ex}

\begin{ex} %[2H9V47]
Hình nào sau đây có thể vừa lăn được vừa trượt được?
\choice
{\True Khối trụ}
{Khối cầu}
{Khối lập phương}
{Khối hộp chữ nhật}
\loigiai{
Khối trụ đặt đứng thì trượt được, đặt nằm thì lăn được.\\
\textbf{Đáp án của câu hỏi là A}
}
\end{ex}

\begin{ex} %[2H9C47]
Có 12 khối lập phương nhỏ. Có thể xếp thành một khối hộp chữ nhật gồm mấy tầng, mỗi tầng 4 khối?
\choice
{$2$ tầng}
{$4$ tầng}
{$5$ tầng}
{\True $3$ tầng}
\loigiai{
$12 : 4 = 3$ tầng.\\
\textbf{Đáp án của câu hỏi là D}
}
\end{ex}

\textbf{Phần 2. Trắc nghiệm đúng sai}
\setcounterpageref{ex}{0}

\begin{ex} %[2H9H47]
Xét tính đúng sai về các hình khối đã học:
\choiceTF
{\True Khối cầu lăn được về mọi phía.}
{\True Khối trụ có 2 mặt phẳng đối diện là hai hình tròn bằng nhau.}
{Khối lập phương có các mặt là hình chữ nhật.}
{\True Khối lập phương có 6 mặt đều là hình vuông bằng nhau.}
\loigiai{
\begin{itemize}
    \item Khối cầu lăn mọi phía: Đúng.
    \item Khối trụ có 2 mặt đáy tròn: Đúng.
    \item Các mặt của khối lập phương là hình vuông. Ý 3 Sai, Ý 4 Đúng.
\end{itemize}
\begin{itemize}
    \item a) \textbf{Đúng}.
    \item b) \textbf{Đúng}.
    \item c) \textbf{Sai}.
    \item d) \textbf{Đúng}.
\end{itemize}
}
\end{ex}

\begin{ex} %[2H9V47]
Bạn Mai có 6 khối cầu và 9 khối trụ. Xét tính đúng sai:
\choiceTF
{\True Tổng số khối hình Mai có là 15 khối ($6 + 9 = 15$).}
{\True Số khối trụ nhiều hơn số khối cầu là 3 khối ($9 - 6 = 3$).}
{Mai có số khối cầu nhiều hơn số khối trụ.}
{\True Mai có số khối cầu ít hơn số khối trụ là 3 khối.}
\loigiai{
\begin{itemize}
    \item Tổng: $6 + 9 = 15$: Đúng.
    \item Trụ hơn cầu: $9 - 6 = 3$: Đúng.
    \item Cầu ít hơn trụ. Ý 3 Sai, Ý 4 Đúng.
\end{itemize}
\begin{itemize}
    \item a) \textbf{Đúng}.
    \item b) \textbf{Đúng}.
    \item c) \textbf{Sai}.
    \item d) \textbf{Đúng}.
\end{itemize}
}
\end{ex}

\textbf{Phần 3. Trả lời ngắn}
\setcounterpageref{ex}{0}

\begin{ex} %[2H9H47]
Khối lập phương có bao nhiêu mặt?
\shortans{6}
\loigiai{
\textbf{Đáp số:} $6$\\
6 mặt.
}
\end{ex}

\begin{ex} %[2H9H47]
Khối trụ có bao nhiêu mặt đáy phẳng hình tròn?
\shortans{2}
\loigiai{
\textbf{Đáp số:} $2$\\
2 mặt đáy.
}
\end{ex}

\begin{ex} %[2H9V47]
Có 7 khối cầu và 8 khối trụ. Hỏi có tất cả bao nhiêu khối hình?
\shortans{15}
\loigiai{
\textbf{Đáp số:} $15$\\
$7 + 8 = 15$ khối.
}
\end{ex}

\begin{ex} %[2H9V47]
Một hộp có 18 khối gỗ gồm khối trụ và khối cầu, trong đó có 9 khối cầu. Hỏi có bao nhiêu khối trụ?
\shortans{9}
\loigiai{
\textbf{Đáp số:} $9$\\
$18 - 9 = 9$ khối.
}
\end{ex}

\begin{ex} %[2H9V47]
Có 20 khối trụ chia đều vào 5 hộp. Hỏi mỗi hộp có bao nhiêu khối trụ?
\shortans{4}
\loigiai{
\textbf{Đáp số:} $4$\\
$20 : 5 = 4$ khối.
}
\end{ex}

\begin{ex} %[2H9C47]
Dùng các khối lập phương nhỏ xếp thành một khối hộp chữ nhật có 3 tầng, mỗi tầng có 6 khối nhỏ. Hỏi cần tất cả bao nhiêu khối lập phương nhỏ?
\shortans{18}
\loigiai{
\textbf{Đáp số:} $18$\\
$6 \times 3 = 18$ khối.
}
\end{ex}
